\documentclass{article}
\usepackage[utf8]{inputenc}
\usepackage[portuguese]{babel}
\usepackage{amsmath}
\usepackage{amsfonts}
\usepackage{amssymb}
\usepackage{geometry}
\geometry{a4paper, margin=1in}

\title{MAC0344 Arquitetura de Computadores \\ Lista de Exercícios No. 4}
\author{Leonardo Heidi Almeida Murakami \\ NUSP: 11260186}
\date{\today}

\begin{document}
\maketitle

\section*{1. Para um dado de 7 bits, escreva o código de Hamming}

\textbf{Enunciado:} Para um dado de 7 bits $m_{1}m_{2}m_{3}m_{4}m_{5}m_{6}m_{7}$, seu código de Hamming tem 11 bits (7 bits originais mais 4 bits adicionais). Escreva o código de Hamming $x_{1}x_{2}x_{3}x_{4}x_{5}x_{6}x_{7}x_{8}x_{9}x_{10}x_{11}$ para o dado $m_{1}m_{2}m_{3}m_{4}m_{5}m_{6}m_{7}=1100101,$ pelo método discutido em classe.

\subsection*{Resolução}

Para começar, precisamos distribuir os 7 bits de dados nas posições que não são potências de 2. Isso acontece porque as posições $x_1, x_2, x_4, x_8$ ficam reservadas para guardar os bits de paridade.

\begin{itemize}
    \item $x_3 \leftarrow m_1 = 1$
    \item $x_5 \leftarrow m_2 = 1$
    \item $x_6 \leftarrow m_3 = 0$
    \item $x_7 \leftarrow m_4 = 0$
    \item $x_9 \leftarrow m_5 = 1$
    \item $x_{10} \leftarrow m_6 = 0$
    \item $x_{11} \leftarrow m_7 = 1$
\end{itemize}

Agora calculamos os bits de paridade ($x_1, x_2, x_4, x_8$). Para isso, aplicamos a operação XOR ($\oplus$) nas posições certas de acordo com o método visto em aula:

\begin{itemize}
    \item \textbf{Cálculo de $x_1$:}
    $x_1 = x_3 \oplus x_5 \oplus x_7 \oplus x_9 \oplus x_{11} = 1 \oplus 1 \oplus 0 \oplus 1 \oplus 1 = 0$
    
    \bigskip
    \item \textbf{Cálculo de $x_2$:}
    $x_2 = x_3 \oplus x_6 \oplus x_7 \oplus x_{10} \oplus x_{11} = 1 \oplus 0 \oplus 0 \oplus 0 \oplus 1 = 0$
    
    \bigskip
    \item \textbf{Cálculo de $x_4$:}
    $x_4 = x_5 \oplus x_6 \oplus x_7 = 1 \oplus 0 \oplus 0 = 1$
    
    \bigskip
    \item \textbf{Cálculo de $x_8$:}
    $x_8 = x_9 \oplus x_{10} \oplus x_{11} = 1 \oplus 0 \oplus 1 = 0$
\end{itemize}

Juntando tudo, o código de Hamming fica:
\begin{center}
\textbf{00111000101}
\end{center}

\newpage
\section*{2. Detecte e corrija o erro no código de Hamming lido}

\textbf{Enunciado:} Suponha que o código de Hamming (total 11 bits) lido da memória foi 
\[y_{1}y_{2}y_{3}y_{4}y_{5}y_{6}y_{7}y_{8}y_{9}y_{10}y_{11}=00110000101.\]
Você detecta algum erro? Se positivo, corrija o erro. Mostre o seu trabalho para chegar a resposta.

\subsection*{Resolução}
Para saber se tem erro e onde ele está, vamos calcular os 4 bits de verificação ($k_1, k_2, k_3, k_4$) com base no código que chegou ($y_i$):

\begin{itemize}
    \item \textbf{Cálculo de $k_1$:}
    $k_1 = y_1 \oplus y_3 \oplus y_5 \oplus y_7 \oplus y_9 \oplus y_{11} = 0 \oplus 1 \oplus 0 \oplus 0 \oplus 1 \oplus 1 = 1$
    
    \bigskip
    \item \textbf{Cálculo de $k_2$:}
    $k_2 = y_2 \oplus y_3 \oplus y_6 \oplus y_7 \oplus y_{10} \oplus y_{11} = 0 \oplus 1 \oplus 0 \oplus 0 \oplus 0 \oplus 1 = 0$
    
    \bigskip
    \item \textbf{Cálculo de $k_3$:}
    $k_3 = y_4 \oplus y_5 \oplus y_6 \oplus y_7 = 1 \oplus 0 \oplus 0 \oplus 0 = 1$
    
    \bigskip
    \item \textbf{Cálculo de $k_4$:}
    $k_4 = y_8 \oplus y_9 \oplus y_{10} \oplus y_{11} = 0 \oplus 1 \oplus 0 \oplus 1 = 0$
\end{itemize}

Juntando os bits de verificação, temos $k_4k_3k_2k_1 = 0101$. Como deu diferente de $0000$, encontramos um erro! E a sacada aqui é que o valor em decimal dessa palavra binária mostra exatamente onde está o erro:
\[
(0101)_2 = 0 \cdot 2^3 + 1 \cdot 2^2 + 0 \cdot 2^1 + 1 \cdot 2^0 = 4 + 1 = 5
\]
ou seja, o problema está na \textbf{posição 5} ($y_5$).

Portanto, flipando o bit da posição 5, o código corrigido fica:
\begin{center}
\textbf{00111000101}
\end{center}

\end{document}